%------------------------------------------------------------------------------%
\begin{question}
  Verifique que a transformação
  \(
    T\left(\vetor{x}{y}\right) = \vetor{2x-y}{x+3y}
  \)
  é linear
\end{question}

%------------------------------------------------------------------------------%
\begin{resolution}{Solução}
  Temos que mostrar que para quais quer vetores
  \[
    \vec{v} = \vetor{r}{s}
    \qquad\qquad
    \vec{w} = \vetor{p}{q}
  \]
  e qualquer escalar $c$

  \[
    T(\vec{v}+\vec{w}) = T(\vec{v}) + T(\vec{w})
  \]
  \[
    T(c\vec{v}) = c T(\vec{v})
  \]
\end{resolution}

%------------------------------------------------------------------------------%
\begin{resolution}{Soma}
  \begin{align*}
    \onslide<+->{T(\vec{v}+\vec{w})}
    \onslide<+->{  = T\left(\vetor{r}{s}+\vetor{p}{q}\right)}
    \onslide<+->{& = T\left(\vetor{r+p}{s+q}\right) \\}
    \onslide<+->{& = \vetor{2x-y}{x+3y} \evalat{x=r+p, \; y=s+q}{} \\}
    \onslide<+->{& = \vetor{2(r+p)-(s+q)}{(r+p)+3(s+q)} \\}
    \onslide<+->{& = \vetor{2p - q + 2r - s}{p + 3q + r + 3s}}
  \end{align*}
\end{resolution}

%------------------------------------------------------------------------------%
\begin{resolution}{Soma}
  \begin{align*}
    \onslide<+->{T(\vec{v}) + T(\vec{w})}
    \onslide<+->{& = T\left(\vetor{r}{s}\right) + T\left(\vetor{p}{q}\right) \\}
    \onslide<+->{& = \vetor{2x-y}{x+3y} \evalat{x=r, \; y=s}{}}
    \onslide<+->{  + \vetor{2x-y}{x+3y} \evalat{x=p, \; y=q}{} \\}
    \onslide<+->{& = \vetor{2r-s}{r+3s} }
    \onslide<+->{  + \vetor{2p-q}{p+3q} \\}
    \onslide<+->{& = \vetor{2p - q + 2r - s}{p + 3q + r + 3s} \\}
    \onslide<+->{& = T(\vec{v}+\vec{w})}
  \end{align*}
\end{resolution}

%------------------------------------------------------------------------------%
\begin{resolution}{Produto por Escalar}
    \begin{align*}
      \onslide<+->{T(c\vec{v})}
      \onslide<+->{  = T\left(\vetor{cr}{cs}\right)}
      \onslide<+->{& = \vetor{2x-y}{x+3y} \evalat{x=cr, \; y=cs}{} \\}
      \onslide<+->{& = \vetor{2cr-cs}{cr+3cs} \\}
      \onslide<+->{& = \vetor{c(2r-s)}{c(r+3s)} \\}
      \onslide<+->{& = c\vetor{2r-s}{r+3s}}
      \onslide<+->{\;=\; c T(\vec{v})}
    \end{align*}
\end{resolution}

%------------------------------------------------------------------------------%
\endinput
\end{frame}
